% 本文件由 tools/build_book.py 从 dist/reports/assertions.json 生成，请勿手改。
\begin{longtable}{@{}p{0.06\textwidth}p{0.32\textwidth}p{0.22\textwidth}p{0.16\textwidth}@{}}
\caption{章节 3.5 数值断言（生成时间 2026-10-19T16:08:50-0400）}\\
\toprule
类别 & 断言 & 理论值 & 实测值 \\
\midrule
\endfirsthead
\toprule
类别 & 断言 & 理论值 & 实测值 \\
\midrule
\endhead
\bottomrule
\endlastfoot

A & 四个周期倍增点与闭式/文献值的最大偏差 & 周期倍增点由乘子穿过 $-$1 定义：r1=\allowbreak{}3(\allowbreak{}不动点失稳)、\allowbreak{}r2=\allowbreak{}1+$\surd$6、\allowbreak{}r3=\allowbreak{}3.5440903595519225、\allowbreak{}r4=\allowbreak{}3.564407266095432(\allowbreak{}后两个是标准文献值) & $1.332\times 10^{-15}$ \checkmark \\
A & r2 与 1+$\surd$6 的偏差 & 周期倍增点由乘子穿过 $-$1 定义：r1=\allowbreak{}3(\allowbreak{}不动点失稳)、\allowbreak{}r2=\allowbreak{}1+$\surd$6、\allowbreak{}r3=\allowbreak{}3.5440903595519225、\allowbreak{}r4=\allowbreak{}3.564407266095432(\allowbreak{}后两个是标准文献值) & 0 \checkmark \\
A & p=1 在 r1 处的乘子(最接近 $-$1 的那条) & 周期倍增点由乘子穿过 $-$1 定义：r1=\allowbreak{}3(\allowbreak{}不动点失稳)、\allowbreak{}r2=\allowbreak{}1+$\surd$6、\allowbreak{}r3=\allowbreak{}3.5440903595519225、\allowbreak{}r4=\allowbreak{}3.564407266095432(\allowbreak{}后两个是标准文献值) & $4.441\times 10^{-16}$ \checkmark \\
A & p=2 在 r2 处的乘子(最接近 $-$1 的那条) & 周期倍增点由乘子穿过 $-$1 定义：r1=\allowbreak{}3(\allowbreak{}不动点失稳)、\allowbreak{}r2=\allowbreak{}1+$\surd$6、\allowbreak{}r3=\allowbreak{}3.5440903595519225、\allowbreak{}r4=\allowbreak{}3.564407266095432(\allowbreak{}后两个是标准文献值) & $6.661\times 10^{-16}$ \checkmark \\
A & p=4 在 r3 处的乘子(最接近 $-$1 的那条) & 周期倍增点由乘子穿过 $-$1 定义：r1=\allowbreak{}3(\allowbreak{}不动点失稳)、\allowbreak{}r2=\allowbreak{}1+$\surd$6、\allowbreak{}r3=\allowbreak{}3.5440903595519225、\allowbreak{}r4=\allowbreak{}3.564407266095432(\allowbreak{}后两个是标准文献值) & $1.754\times 10^{-14}$ \checkmark \\
A & p=8 在 r4 处的乘子(最接近 $-$1 的那条) & 周期倍增点由乘子穿过 $-$1 定义：r1=\allowbreak{}3(\allowbreak{}不动点失稳)、\allowbreak{}r2=\allowbreak{}1+$\surd$6、\allowbreak{}r3=\allowbreak{}3.5440903595519225、\allowbreak{}r4=\allowbreak{}3.564407266095432(\allowbreak{}后两个是标准文献值) & $1.765\times 10^{-14}$ \checkmark \\
A & r=3.2 的 2 周期点 vs 闭式 & logistic 的 2 周期有闭式：x=\allowbreak{}(\allowbreak{}r+1$\pm$$\surd$(\allowbreak{}(\allowbreak{}r$-$3)(\allowbreak{}r+1)))/\allowbreak{}(\allowbreak{}2r)，乘子 M=\allowbreak{}g'(\allowbreak{}x$_1$)g'(\allowbreak{}x$_2$)=\allowbreak{}4+2r$-$r$^2$ & $1.110\times 10^{-16}$ \checkmark \\
A & r=3.2 的乘子 vs 4+2r$-$r$^2$ & logistic 的 2 周期有闭式：x=\allowbreak{}(\allowbreak{}r+1$\pm$$\surd$(\allowbreak{}(\allowbreak{}r$-$3)(\allowbreak{}r+1)))/\allowbreak{}(\allowbreak{}2r)，乘子 M=\allowbreak{}g'(\allowbreak{}x$_1$)g'(\allowbreak{}x$_2$)=\allowbreak{}4+2r$-$r$^2$ & $6.384\times 10^{-16}$ \checkmark \\
A & r=3.2 的两条乘子路径之差(Newton vs 网格二分) & logistic 的 2 周期有闭式：x=\allowbreak{}(\allowbreak{}r+1$\pm$$\surd$(\allowbreak{}(\allowbreak{}r$-$3)(\allowbreak{}r+1)))/\allowbreak{}(\allowbreak{}2r)，乘子 M=\allowbreak{}g'(\allowbreak{}x$_1$)g'(\allowbreak{}x$_2$)=\allowbreak{}4+2r$-$r$^2$ & $3.469\times 10^{-15}$ \checkmark \\
A & r=3.4 的 2 周期点 vs 闭式 & logistic 的 2 周期有闭式：x=\allowbreak{}(\allowbreak{}r+1$\pm$$\surd$(\allowbreak{}(\allowbreak{}r$-$3)(\allowbreak{}r+1)))/\allowbreak{}(\allowbreak{}2r)，乘子 M=\allowbreak{}g'(\allowbreak{}x$_1$)g'(\allowbreak{}x$_2$)=\allowbreak{}4+2r$-$r$^2$ & $2.220\times 10^{-16}$ \checkmark \\
A & r=3.4 的乘子 vs 4+2r$-$r$^2$ & logistic 的 2 周期有闭式：x=\allowbreak{}(\allowbreak{}r+1$\pm$$\surd$(\allowbreak{}(\allowbreak{}r$-$3)(\allowbreak{}r+1)))/\allowbreak{}(\allowbreak{}2r)，乘子 M=\allowbreak{}g'(\allowbreak{}x$_1$)g'(\allowbreak{}x$_2$)=\allowbreak{}4+2r$-$r$^2$ & $1.332\times 10^{-15}$ \checkmark \\
A & r=3.4 的两条乘子路径之差(Newton vs 网格二分) & logistic 的 2 周期有闭式：x=\allowbreak{}(\allowbreak{}r+1$\pm$$\surd$(\allowbreak{}(\allowbreak{}r$-$3)(\allowbreak{}r+1)))/\allowbreak{}(\allowbreak{}2r)，乘子 M=\allowbreak{}g'(\allowbreak{}x$_1$)g'(\allowbreak{}x$_2$)=\allowbreak{}4+2r$-$r$^2$ & $1.554\times 10^{-15}$ \checkmark \\
A & r=3.23607 的 2 周期点 vs 闭式 & logistic 的 2 周期有闭式：x=\allowbreak{}(\allowbreak{}r+1$\pm$$\surd$(\allowbreak{}(\allowbreak{}r$-$3)(\allowbreak{}r+1)))/\allowbreak{}(\allowbreak{}2r)，乘子 M=\allowbreak{}g'(\allowbreak{}x$_1$)g'(\allowbreak{}x$_2$)=\allowbreak{}4+2r$-$r$^2$ & $1.110\times 10^{-16}$ \checkmark \\
A & r=3.23607 的乘子 vs 4+2r$-$r$^2$ & logistic 的 2 周期有闭式：x=\allowbreak{}(\allowbreak{}r+1$\pm$$\surd$(\allowbreak{}(\allowbreak{}r$-$3)(\allowbreak{}r+1)))/\allowbreak{}(\allowbreak{}2r)，乘子 M=\allowbreak{}g'(\allowbreak{}x$_1$)g'(\allowbreak{}x$_2$)=\allowbreak{}4+2r$-$r$^2$ & $1.437\times 10^{-15}$ \checkmark \\
A & r=3.23607 的两条乘子路径之差(Newton vs 网格二分) & logistic 的 2 周期有闭式：x=\allowbreak{}(\allowbreak{}r+1$\pm$$\surd$(\allowbreak{}(\allowbreak{}r$-$3)(\allowbreak{}r+1)))/\allowbreak{}(\allowbreak{}2r)，乘子 M=\allowbreak{}g'(\allowbreak{}x$_1$)g'(\allowbreak{}x$_2$)=\allowbreak{}4+2r$-$r$^2$ & $2.874\times 10^{-15}$ \checkmark \\
A & 3 周期窗口起点 vs 1+$\surd$8 & 3 周期窗口起点是切分岔：g$^3$(\allowbreak{}x)$-$x 出现二重根，闭式 r=\allowbreak{}1+$\surd$8=\allowbreak{}3.8284271247461903 & $4.685\times 10^{-10}$ \checkmark \\
A & 窗口内外的乘子(内 |M|<1、外无 3 周期) & 3 周期窗口起点是切分岔：g$^3$(\allowbreak{}x)$-$x 出现二重根，闭式 r=\allowbreak{}1+$\surd$8=\allowbreak{}3.8284271247461903 & 0.83293 \checkmark \\
C & 前六个超稳定参数与文献值的最大偏差 & Feigenbaum 常数 $\delta$=\allowbreak{}lim (\allowbreak{}R\_n$-$R\_\{n$-$1\})/\allowbreak{}(\allowbreak{}R\_\{n+1\}$-$R\_n)=\allowbreak{}4.6692016091$\dots$，其中 R\_n 是超稳定参数(\allowbreak{}2\textasciicircum{}\{n$-$1\} 周期轨经过 x=\allowbreak{}1/\allowbreak{}2)。序列收敛不快也不单调，所以容差按实测尾项给 & $4.683\times 10^{-11}$ \checkmark \\
C & R$_2$ 与 1+$\surd$5 的偏差 & Feigenbaum 常数 $\delta$=\allowbreak{}lim (\allowbreak{}R\_n$-$R\_\{n$-$1\})/\allowbreak{}(\allowbreak{}R\_\{n+1\}$-$R\_n)=\allowbreak{}4.6692016091$\dots$，其中 R\_n 是超稳定参数(\allowbreak{}2\textasciicircum{}\{n$-$1\} 周期轨经过 x=\allowbreak{}1/\allowbreak{}2)。序列收敛不快也不单调，所以容差按实测尾项给 & 0 \checkmark \\
C & $\delta$ 序列尾项与 4.6692016091 之差 & Feigenbaum 常数 $\delta$=\allowbreak{}lim (\allowbreak{}R\_n$-$R\_\{n$-$1\})/\allowbreak{}(\allowbreak{}R\_\{n+1\}$-$R\_n)=\allowbreak{}4.6692016091$\dots$，其中 R\_n 是超稳定参数(\allowbreak{}2\textasciicircum{}\{n$-$1\} 周期轨经过 x=\allowbreak{}1/\allowbreak{}2)。序列收敛不快也不单调，所以容差按实测尾项给 & $1.061\times 10^{-5}$ \checkmark \\
C & $\delta$ 序列是否在收敛(尾项误差 < 中段误差) & Feigenbaum 常数 $\delta$=\allowbreak{}lim (\allowbreak{}R\_n$-$R\_\{n$-$1\})/\allowbreak{}(\allowbreak{}R\_\{n+1\}$-$R\_n)=\allowbreak{}4.6692016091$\dots$，其中 R\_n 是超稳定参数(\allowbreak{}2\textasciicircum{}\{n$-$1\} 周期轨经过 x=\allowbreak{}1/\allowbreak{}2)。序列收敛不快也不单调，所以容差按实测尾项给 & $1.061\times 10^{-5}$ \checkmark \\
C & 外推的聚点 r$\infty$ 与文献值 3.5699456718 之差 & Feigenbaum 常数 $\delta$=\allowbreak{}lim (\allowbreak{}R\_n$-$R\_\{n$-$1\})/\allowbreak{}(\allowbreak{}R\_\{n+1\}$-$R\_n)=\allowbreak{}4.6692016091$\dots$，其中 R\_n 是超稳定参数(\allowbreak{}2\textasciicircum{}\{n$-$1\} 周期轨经过 x=\allowbreak{}1/\allowbreak{}2)。序列收敛不快也不单调，所以容差按实测尾项给 & $6.763\times 10^{-11}$ \checkmark \\
A & r=3.2 的最小周期 & logistic 在 r=\allowbreak{}3.2 是 2 周期、\allowbreak{}r=\allowbreak{}3.5 是 4 周期、\allowbreak{}r=\allowbreak{}3.83/\allowbreak{}3.84 落在 3 周期窗口里、\allowbreak{}r=\allowbreak{}4 是混沌(\allowbreak{}$\le$32 步内不闭合)。判据是 |f\textasciicircum{}p(\allowbreak{}x)$-$x|$\le$1e-9，不是肉眼 & 2 \checkmark \\
A & r=3.5 的最小周期 & logistic 在 r=\allowbreak{}3.2 是 2 周期、\allowbreak{}r=\allowbreak{}3.5 是 4 周期、\allowbreak{}r=\allowbreak{}3.83/\allowbreak{}3.84 落在 3 周期窗口里、\allowbreak{}r=\allowbreak{}4 是混沌(\allowbreak{}$\le$32 步内不闭合)。判据是 |f\textasciicircum{}p(\allowbreak{}x)$-$x|$\le$1e-9，不是肉眼 & 4 \checkmark \\
A & r=3.83 的最小周期 & logistic 在 r=\allowbreak{}3.2 是 2 周期、\allowbreak{}r=\allowbreak{}3.5 是 4 周期、\allowbreak{}r=\allowbreak{}3.83/\allowbreak{}3.84 落在 3 周期窗口里、\allowbreak{}r=\allowbreak{}4 是混沌(\allowbreak{}$\le$32 步内不闭合)。判据是 |f\textasciicircum{}p(\allowbreak{}x)$-$x|$\le$1e-9，不是肉眼 & 3 \checkmark \\
A & r=3.84 的最小周期 & logistic 在 r=\allowbreak{}3.2 是 2 周期、\allowbreak{}r=\allowbreak{}3.5 是 4 周期、\allowbreak{}r=\allowbreak{}3.83/\allowbreak{}3.84 落在 3 周期窗口里、\allowbreak{}r=\allowbreak{}4 是混沌(\allowbreak{}$\le$32 步内不闭合)。判据是 |f\textasciicircum{}p(\allowbreak{}x)$-$x|$\le$1e-9，不是肉眼 & 3 \checkmark \\
A & r=4 的周期(应为 0 = 32 步内不闭合) & logistic 在 r=\allowbreak{}3.2 是 2 周期、\allowbreak{}r=\allowbreak{}3.5 是 4 周期、\allowbreak{}r=\allowbreak{}3.83/\allowbreak{}3.84 落在 3 周期窗口里、\allowbreak{}r=\allowbreak{}4 是混沌(\allowbreak{}$\le$32 步内不闭合)。判据是 |f\textasciicircum{}p(\allowbreak{}x)$-$x|$\le$1e-9，不是肉眼 & 0 \checkmark \\
A & 周期窗口里的最大闭合残差 & 周期窗口里闭合残差是舍入水平(\allowbreak{}1e-15)，混沌区是 O(\allowbreak{}1e-2)：两类点必须分在对数轴两端；混沌区的残差**不能**是 0(\allowbreak{}那会把混沌画成完美闭合) & $9.279\times 10^{-10}$ \checkmark \\
A & 混沌区闭合残差的中位数 & 周期窗口里闭合残差是舍入水平(\allowbreak{}1e-15)，混沌区是 O(\allowbreak{}1e-2)：两类点必须分在对数轴两端；混沌区的残差**不能**是 0(\allowbreak{}那会把混沌画成完美闭合) & 0.005202 \checkmark \\
A & 混沌区里残差 > 1e-4 的参数占比 & 周期窗口里闭合残差是舍入水平(\allowbreak{}1e-15)，混沌区是 O(\allowbreak{}1e-2)：两类点必须分在对数轴两端；混沌区的残差**不能**是 0(\allowbreak{}那会把混沌画成完美闭合) & 0.943503 \checkmark \\
A & 向量化 vs 标量的周期差异计数 & 向量化的 detect\_periods 与逐参数标量 detect\_period 必须给出同一个周期与同一个残差(\allowbreak{}批量路径只是把参数维摊开；残差允许舍入量级的差) & 0 \checkmark \\
A & 向量化 vs 标量的残差最大差 & 向量化的 detect\_periods 与逐参数标量 detect\_period 必须给出同一个周期与同一个残差(\allowbreak{}批量路径只是把参数维摊开；残差允许舍入量级的差) & $3.220\times 10^{-15}$ \checkmark \\
A & r=2.5 的不动点是否稳定 & r=\allowbreak{}4 的超稳定不动点(\allowbreak{}x=\allowbreak{}1/\allowbreak{}2 是 r=\allowbreak{}2 的临界点)与周期窗口的存在性给出 Logistic 映射分岔结构的两个锚点：r<3 只有 1 周期、\allowbreak{}r=\allowbreak{}3.2 的乘子为 0.16 & 1 \checkmark \\
A & r=3.2 的不动点是否稳定 & r=\allowbreak{}4 的超稳定不动点(\allowbreak{}x=\allowbreak{}1/\allowbreak{}2 是 r=\allowbreak{}2 的临界点)与周期窗口的存在性给出 Logistic 映射分岔结构的两个锚点：r<3 只有 1 周期、\allowbreak{}r=\allowbreak{}3.2 的乘子为 0.16 & 0 \checkmark \\
A & r=3.2 处不动点的乘子里有没有 2$-$r=$-$1.2 & r=\allowbreak{}4 的超稳定不动点(\allowbreak{}x=\allowbreak{}1/\allowbreak{}2 是 r=\allowbreak{}2 的临界点)与周期窗口的存在性给出 Logistic 映射分岔结构的两个锚点：r<3 只有 1 周期、\allowbreak{}r=\allowbreak{}3.2 的乘子为 0.16 & 0 \checkmark \\
A & r=3.2 处稳定 2 周期的乘子 vs 4+2r$-$r$^2$=0.16 & r=\allowbreak{}4 的超稳定不动点(\allowbreak{}x=\allowbreak{}1/\allowbreak{}2 是 r=\allowbreak{}2 的临界点)与周期窗口的存在性给出 Logistic 映射分岔结构的两个锚点：r<3 只有 1 周期、\allowbreak{}r=\allowbreak{}3.2 的乘子为 0.16 & $2.276\times 10^{-15}$ \checkmark \\
A & 三个平衡点的 |f| 残差 & Lorenz 的非平凡平衡点 C$\pm$=\allowbreak{}(\allowbreak{}$\pm$$\surd$(\allowbreak{}$\beta$(\allowbreak{}$\rho$$-$1)),\allowbreak{} $\pm$$\surd$(\allowbreak{}$\beta$(\allowbreak{}$\rho$$-$1)),\allowbreak{} $\rho$$-$1) 与原点都是 f=\allowbreak{}0 的精确解；散度 $\nabla$$\cdot$f=\allowbreak{}$-$(\allowbreak{}$\sigma$+1+$\beta$) 与状态无关(\allowbreak{}常数) & $1.421\times 10^{-14}$ \checkmark \\
A & C$\pm$ 与闭式 $\surd$($\beta$($\rho$$-$1)) 的偏差 & Lorenz 的非平凡平衡点 C$\pm$=\allowbreak{}(\allowbreak{}$\pm$$\surd$(\allowbreak{}$\beta$(\allowbreak{}$\rho$$-$1)),\allowbreak{} $\pm$$\surd$(\allowbreak{}$\beta$(\allowbreak{}$\rho$$-$1)),\allowbreak{} $\rho$$-$1) 与原点都是 f=\allowbreak{}0 的精确解；散度 $\nabla$$\cdot$f=\allowbreak{}$-$(\allowbreak{}$\sigma$+1+$\beta$) 与状态无关(\allowbreak{}常数) & 0 \checkmark \\
A & 散度 = $-$($\sigma$+1+$\beta$) & Lorenz 的非平凡平衡点 C$\pm$=\allowbreak{}(\allowbreak{}$\pm$$\surd$(\allowbreak{}$\beta$(\allowbreak{}$\rho$$-$1)),\allowbreak{} $\pm$$\surd$(\allowbreak{}$\beta$(\allowbreak{}$\rho$$-$1)),\allowbreak{} $\rho$$-$1) 与原点都是 f=\allowbreak{}0 的精确解；散度 $\nabla$$\cdot$f=\allowbreak{}$-$(\allowbreak{}$\sigma$+1+$\beta$) 与状态无关(\allowbreak{}常数) & 0 \checkmark \\
A & 原点的三个实特征值 & Lorenz 的非平凡平衡点 C$\pm$=\allowbreak{}(\allowbreak{}$\pm$$\surd$(\allowbreak{}$\beta$(\allowbreak{}$\rho$$-$1)),\allowbreak{} $\pm$$\surd$(\allowbreak{}$\beta$(\allowbreak{}$\rho$$-$1)),\allowbreak{} $\rho$$-$1) 与原点都是 f=\allowbreak{}0 的精确解；散度 $\nabla$$\cdot$f=\allowbreak{}$-$(\allowbreak{}$\sigma$+1+$\beta$) 与状态无关(\allowbreak{}常数) & $5.488\times 10^{-7}$ \checkmark \\
A & C$\pm$ 复共轭对的虚部与 10.194501 之差 & Lorenz 的非平凡平衡点 C$\pm$=\allowbreak{}(\allowbreak{}$\pm$$\surd$(\allowbreak{}$\beta$(\allowbreak{}$\rho$$-$1)),\allowbreak{} $\pm$$\surd$(\allowbreak{}$\beta$(\allowbreak{}$\rho$$-$1)),\allowbreak{} $\rho$$-$1) 与原点都是 f=\allowbreak{}0 的精确解；散度 $\nabla$$\cdot$f=\allowbreak{}$-$(\allowbreak{}$\sigma$+1+$\beta$) 与状态无关(\allowbreak{}常数) & $4.221\times 10^{-6}$ \checkmark \\
A & $\rho$=28 处 C$\pm$ 的相型(鞍焦) & Lorenz 的非平凡平衡点 C$\pm$=\allowbreak{}(\allowbreak{}$\pm$$\surd$(\allowbreak{}$\beta$(\allowbreak{}$\rho$$-$1)),\allowbreak{} $\pm$$\surd$(\allowbreak{}$\beta$(\allowbreak{}$\rho$$-$1)),\allowbreak{} $\rho$$-$1) 与原点都是 f=\allowbreak{}0 的精确解；散度 $\nabla$$\cdot$f=\allowbreak{}$-$(\allowbreak{}$\sigma$+1+$\beta$) 与状态无关(\allowbreak{}常数) & 7 \checkmark \\
C & 二分实测 $\rho$\_\allowbreak{}H 与 470/19 & Hopf 分岔的定义是平衡点的一对特征值穿过虚轴：max Re $\lambda$(\allowbreak{}$\rho$) 在 $\rho$\_H 变号。二分实测该变号点，与闭式 $\rho$\_H=\allowbreak{}$\sigma$(\allowbreak{}$\sigma$+$\beta$+3)/\allowbreak{}(\allowbreak{}$\sigma$$-$$\beta$$-$1)=\allowbreak{}470/\allowbreak{}19 对拍 & $9.809\times 10^{-12}$ \checkmark \\
C & $\rho$\_\allowbreak{}H 处 C$\pm$ 的 max Re $\lambda$(应为 0) & Hopf 分岔的定义是平衡点的一对特征值穿过虚轴：max Re $\lambda$(\allowbreak{}$\rho$) 在 $\rho$\_H 变号。二分实测该变号点，与闭式 $\rho$\_H=\allowbreak{}$\sigma$(\allowbreak{}$\sigma$+$\beta$+3)/\allowbreak{}(\allowbreak{}$\sigma$$-$$\beta$$-$1)=\allowbreak{}470/\allowbreak{}19 对拍 & $5.928\times 10^{-16}$ \checkmark \\
C & $\rho$\_\allowbreak{}H 处第三个特征值(负实根，与 $-$($\sigma$+1+$\beta$) 同值) & Hopf 分岔的定义是平衡点的一对特征值穿过虚轴：max Re $\lambda$(\allowbreak{}$\rho$) 在 $\rho$\_H 变号。二分实测该变号点，与闭式 $\rho$\_H=\allowbreak{}$\sigma$(\allowbreak{}$\sigma$+$\beta$+3)/\allowbreak{}(\allowbreak{}$\sigma$$-$$\beta$$-$1)=\allowbreak{}470/\allowbreak{}19 对拍 & $1.776\times 10^{-15}$ \checkmark \\
A & van der Pol 的 dy/dt=0 零倾线分支数(数值 vs 闭式) & van der Pol 的 dy/\allowbreak{}dt=\allowbreak{}0 零倾线有闭式 y=\allowbreak{}x/\allowbreak{}(\allowbreak{}$\mu$(\allowbreak{}1$-$x$^2$))；数值追踪(\allowbreak{}逐列二分 + 归一化二维距离连接)必须给出同一条曲线，且分支数一致 & 3 \checkmark \\
A & 数值零倾线 vs 闭式 y=x/($\mu$(1$-$x$^2$)) 的最大偏差 & van der Pol 的 dy/\allowbreak{}dt=\allowbreak{}0 零倾线有闭式 y=\allowbreak{}x/\allowbreak{}(\allowbreak{}$\mu$(\allowbreak{}1$-$x$^2$))；数值追踪(\allowbreak{}逐列二分 + 归一化二维距离连接)必须给出同一条曲线，且分支数一致 & $4.441\times 10^{-16}$ \checkmark \\
A & 跨过 x=$-$1 渐近线的分支数(应为 0) & van der Pol 的 dy/\allowbreak{}dt=\allowbreak{}0 零倾线有闭式 y=\allowbreak{}x/\allowbreak{}(\allowbreak{}$\mu$(\allowbreak{}1$-$x$^2$))；数值追踪(\allowbreak{}逐列二分 + 归一化二维距离连接)必须给出同一条曲线，且分支数一致 & 0 \checkmark \\
A & 阻尼摆 dy/dt=0 零倾线 vs 闭式 y=$-$sin(x)/c & van der Pol 的 dy/\allowbreak{}dt=\allowbreak{}0 零倾线有闭式 y=\allowbreak{}x/\allowbreak{}(\allowbreak{}$\mu$(\allowbreak{}1$-$x$^2$))；数值追踪(\allowbreak{}逐列二分 + 归一化二维距离连接)必须给出同一条曲线，且分支数一致 & $4.441\times 10^{-16}$ \checkmark \\
A & 截面点离平面的最大距离(精化后) & 定向 Poincare 截面只收单向穿越：截面上每一点都必须满足 f(\allowbreak{}y)[平面分量]>0(\allowbreak{}向上穿越)且落在平面上；返回时间与穿越顺序一一对应 & $7.280\times 10^{-12}$ \checkmark \\
A & 穿越点上的 dz/dt 全为正(定向判据) & 定向 Poincare 截面只收单向穿越：截面上每一点都必须满足 f(\allowbreak{}y)[平面分量]>0(\allowbreak{}向上穿越)且落在平面上；返回时间与穿越顺序一一对应 & 69.0878 \checkmark \\
A & 相邻穿越的时间间隔(均值$\pm$极差) & 定向 Poincare 截面只收单向穿越：截面上每一点都必须满足 f(\allowbreak{}y)[平面分量]>0(\allowbreak{}向上穿越)且落在平面上；返回时间与穿越顺序一一对应 & 0.748659 \checkmark \\
C & 精化位移随步长的收敛(最大位移) & 截面点的定位误差有两个来源：线性插值是 O(\allowbreak{}h$^2$)，Newton 精化把它压到舍入水平。「线性插值点 $\to$ 精化点」的位移必须按 h$^2$ 收敛 & 0.0153331 \checkmark \\
C & 线性插值定位误差的收敛阶 & 2 & 1.88789 \checkmark \\
C & 平面残差(Newton 精化后) & 截面点的定位误差有两个来源：线性插值是 O(\allowbreak{}h$^2$)，Newton 精化把它压到舍入水平。「线性插值点 $\to$ 精化点」的位移必须按 h$^2$ 收敛 & $1.421\times 10^{-14}$ \checkmark \\
A & $\mu$=0 的周期 $-$ 2$\pi$ & $\mu$=\allowbreak{}0 的 van der Pol 退化成简谐振子 x''+x=\allowbreak{}0，周期精确等于 2$\pi$；$\mu$$\to$0 时极限环周期的修正是 (\allowbreak{}T$-$2$\pi$)/\allowbreak{}$\mu$$^2$ $\to$ 2$\pi$/\allowbreak{}16 & $5.243\times 10^{-10}$ \checkmark \\
A & (T$-$2$\pi$)/$\mu$$^2$ 在 $\mu$=0.02 处的值 vs 2$\pi$/16 & $\mu$=\allowbreak{}0 的 van der Pol 退化成简谐振子 x''+x=\allowbreak{}0，周期精确等于 2$\pi$；$\mu$$\to$0 时极限环周期的修正是 (\allowbreak{}T$-$2$\pi$)/\allowbreak{}$\mu$$^2$ $\to$ 2$\pi$/\allowbreak{}16 & $6.716\times 10^{-6}$ \checkmark \\
A & 极限环振幅随 $\mu$ 的变化(2 的修正) & $\mu$=\allowbreak{}0 的 van der Pol 退化成简谐振子 x''+x=\allowbreak{}0，周期精确等于 2$\pi$；$\mu$$\to$0 时极限环周期的修正是 (\allowbreak{}T$-$2$\pi$)/\allowbreak{}$\mu$$^2$ $\to$ 2$\pi$/\allowbreak{}16 & 0.00861986 \checkmark \\
A & 极限环半径 vs $\surd$$\mu$ & Hopf 标准型 $\rho$'=\allowbreak{}$\mu$$\rho$$-$$\rho$$^3$、\allowbreak{}$\theta$'=\allowbreak{}$\omega$ 的极限环是半径 $\surd$$\mu$ 的圆，周期精确等于 2$\pi$/\allowbreak{}$\omega$ & $2.220\times 10^{-16}$ \checkmark \\
A & 极限环周期 vs 2$\pi$/$\omega$ & Hopf 标准型 $\rho$'=\allowbreak{}$\mu$$\rho$$-$$\rho$$^3$、\allowbreak{}$\theta$'=\allowbreak{}$\omega$ 的极限环是半径 $\surd$$\mu$ 的圆，周期精确等于 2$\pi$/\allowbreak{}$\omega$ & $8.517\times 10^{-12}$ \checkmark \\
A & 批量积分 vs 单条积分(逐点最大差) & integrate\_ensemble 把 (\allowbreak{}n,\allowbreak{}dim) 摊平后用同一个 rk\_step 推进；n=\allowbreak{}1 时它必须与单条轨迹的 integrate\_fixed 逐点相同(\allowbreak{}两条独立的代码路径) & 0 \checkmark \\
A & lorenz\_\allowbreak{}butterfly：由 x/y/z + strings 角度重算 u/v 的偏差 & 三维帧必须同时给原始 x/\allowbreak{}y/\allowbreak{}z 与投影后的 u/\allowbreak{}v，并把投影参数写进 strings；帧里的 u/\allowbreak{}v 必须能由同帧的 x/\allowbreak{}y/\allowbreak{}z 与同一组角度重算出来 & 0 \checkmark \\
A & poincare\_\allowbreak{}section：由 x/y/z + strings 角度重算 u/v 的偏差 & 三维帧必须同时给原始 x/\allowbreak{}y/\allowbreak{}z 与投影后的 u/\allowbreak{}v，并把投影参数写进 strings；帧里的 u/\allowbreak{}v 必须能由同帧的 x/\allowbreak{}y/\allowbreak{}z 与同一组角度重算出来 & 0 \checkmark \\
A & meta.t\_\allowbreak{}visible vs 数组里第一次越过 1e-3 的时刻 & 帧里报出来的标量必须与同帧的数组一致(\allowbreak{}t\_visible 是 sep\_mean 第一次达到显示分辨率１e-3 的时刻)------否则前端画的与 meta 说的是两回事 & 0 \checkmark \\
A & 孪生实验的拟合增长率(默认参数) & 帧里报出来的标量必须与同帧的数组一致(\allowbreak{}t\_visible 是 sep\_mean 第一次达到显示分辨率１e-3 的时刻)------否则前端画的与 meta 说的是两回事 & 0.887642 \checkmark \\
A & h 与 2h 两条分离曲线的增长率之差(默认参数) & 帧里报出来的标量必须与同帧的数组一致(\allowbreak{}t\_visible 是 sep\_mean 第一次达到显示分辨率１e-3 的时刻)------否则前端画的与 meta 说的是两回事 & 0.00195081 \checkmark \\
A & 第 1 条初值的 $\Sigma$$\lambda$ 与散度之差 & Lyapunov 指数的和等于散度的时间平均：$\Sigma$$\lambda$\_i =\allowbreak{} $\langle$$\nabla$$\cdot$f$\rangle$。Lorenz 的散度是常数 $-$(\allowbreak{}$\sigma$+1+$\beta$)=\allowbreak{}$-$13.6667，所以这条恒等式给出一个不依赖任何文献值的标尺 & 0.000102712 \checkmark \\
A & 第 1 条初值的 $\lambda$$_2$(流向，应趋于 0) & Lyapunov 指数的和等于散度的时间平均：$\Sigma$$\lambda$\_i =\allowbreak{} $\langle$$\nabla$$\cdot$f$\rangle$。Lorenz 的散度是常数 $-$(\allowbreak{}$\sigma$+1+$\beta$)=\allowbreak{}$-$13.6667，所以这条恒等式给出一个不依赖任何文献值的标尺 & 0.000324632 \checkmark \\
A & 第 2 条初值的 $\Sigma$$\lambda$ 与散度之差 & Lyapunov 指数的和等于散度的时间平均：$\Sigma$$\lambda$\_i =\allowbreak{} $\langle$$\nabla$$\cdot$f$\rangle$。Lorenz 的散度是常数 $-$(\allowbreak{}$\sigma$+1+$\beta$)=\allowbreak{}$-$13.6667，所以这条恒等式给出一个不依赖任何文献值的标尺 & 0.000101611 \checkmark \\
A & 第 2 条初值的 $\lambda$$_2$(流向，应趋于 0) & Lyapunov 指数的和等于散度的时间平均：$\Sigma$$\lambda$\_i =\allowbreak{} $\langle$$\nabla$$\cdot$f$\rangle$。Lorenz 的散度是常数 $-$(\allowbreak{}$\sigma$+1+$\beta$)=\allowbreak{}$-$13.6667，所以这条恒等式给出一个不依赖任何文献值的标尺 & 0.00062247 \checkmark \\
A & Gram--Schmidt 的最大正交性缺陷 & Lyapunov 指数的和等于散度的时间平均：$\Sigma$$\lambda$\_i =\allowbreak{} $\langle$$\nabla$$\cdot$f$\rangle$。Lorenz 的散度是常数 $-$(\allowbreak{}$\sigma$+1+$\beta$)=\allowbreak{}$-$13.6667，所以这条恒等式给出一个不依赖任何文献值的标尺 & $6.661\times 10^{-16}$ \checkmark \\
A & Henon 的 $\Sigma$$\lambda$ 与 log|b| 之差 & Henon 映射的 Jacobi 行列式处处等于 $-$b，所以两个 Lyapunov 指数之和精确等于 log|b|(\allowbreak{}b=\allowbreak{}0.3 $\to$ $-$1.2039728043259361)，与轨道无关 & $1.177\times 10^{-14}$ \checkmark \\
A & $\lambda$$_1$ 与文献值 0.41922 之差 & Henon 映射的 Jacobi 行列式处处等于 $-$b，所以两个 Lyapunov 指数之和精确等于 log|b|(\allowbreak{}b=\allowbreak{}0.3 $\to$ $-$1.2039728043259361)，与轨道无关 & 0.00038919 \checkmark \\
A & $\lambda$(2.5) 与闭式之差 & logistic 映射的 Lyapunov 指数有三处闭式：1<r<3 时 $\lambda$=\allowbreak{}log|2$-$r|；2 周期窗口内 $\lambda$=\allowbreak{}$1/\allowbreak{}2$log|4+2r$-$r$^2$|；r=\allowbreak{}4 共轭到 Bernoulli 移位，$\lambda$=\allowbreak{}log2 & $3.819\times 10^{-14}$ \checkmark \\
A & $\lambda$(3.2) 与闭式之差 & logistic 映射的 Lyapunov 指数有三处闭式：1<r<3 时 $\lambda$=\allowbreak{}log|2$-$r|；2 周期窗口内 $\lambda$=\allowbreak{}$1/\allowbreak{}2$log|4+2r$-$r$^2$|；r=\allowbreak{}4 共轭到 Bernoulli 移位，$\lambda$=\allowbreak{}log2 & $1.088\times 10^{-14}$ \checkmark \\
A & $\lambda$(4.0) 与闭式之差 & logistic 映射的 Lyapunov 指数有三处闭式：1<r<3 时 $\lambda$=\allowbreak{}log|2$-$r|；2 周期窗口内 $\lambda$=\allowbreak{}$1/\allowbreak{}2$log|4+2r$-$r$^2$|；r=\allowbreak{}4 共轭到 Bernoulli 移位，$\lambda$=\allowbreak{}log2 & $1.847\times 10^{-6}$ \checkmark \\
A & $\lambda$$_1$ 的均值与公认值 0.9056 之差 & Lorenz(\allowbreak{}$\sigma$=\allowbreak{}10,\allowbreak{}$\rho$=\allowbreak{}28,\allowbreak{}$\beta$=\allowbreak{}8/\allowbreak{}3) 的最大 Lyapunov 指数公认值 0.9056(\allowbreak{}Sprott 2003 表 5.1)。有限时间估计 T=\allowbreak{}2000 的涨落由跨初值的标准差实测，容差据此给出 & 0.00154537 \checkmark \\
A & 公认值落在实测的 $\pm$2$\sigma$ 区间内 & Lorenz(\allowbreak{}$\sigma$=\allowbreak{}10,\allowbreak{}$\rho$=\allowbreak{}28,\allowbreak{}$\beta$=\allowbreak{}8/\allowbreak{}3) 的最大 Lyapunov 指数公认值 0.9056(\allowbreak{}Sprott 2003 表 5.1)。有限时间估计 T=\allowbreak{}2000 的涨落由跨初值的标准差实测，容差据此给出 & 0.150413 \checkmark \\
C & h 减半后 $\lambda$$_1$ 的变化 & 定步长积分的 Lyapunov 估计必须随步长收敛：h 减半后 $\lambda$$_1$ 的变化量要小于积分误差允许的量级，否则「混沌」可能只是离散化噪声 & 0.0143774 \checkmark \\
C & h=0.01 的 $\Sigma$$\lambda$ 与散度之差 & 定步长积分的 Lyapunov 估计必须随步长收敛：h 减半后 $\lambda$$_1$ 的变化量要小于积分误差允许的量级，否则「混沌」可能只是离散化噪声 & 0.000103136 \checkmark \\
C & h=0.005 的 $\Sigma$$\lambda$ 与散度之差 & 定步长积分的 Lyapunov 估计必须随步长收敛：h 减半后 $\lambda$$_1$ 的变化量要小于积分误差允许的量级，否则「混沌」可能只是离散化噪声 & $6.215\times 10^{-6}$ \checkmark \\
C & 孪生分离曲线的增长率(h 与 2h)之差 & 定步长积分的 Lyapunov 估计必须随步长收敛：h 减半后 $\lambda$$_1$ 的变化量要小于积分误差允许的量级，否则「混沌」可能只是离散化噪声 & 0.00277876 \checkmark \\
C & QR 间隔 10 vs 20 的谱差 & QR 重正交化的间隔只影响浮点误差与代价，不影响数学：间隔 10 与 20 必须给出同一个谱 & $2.132\times 10^{-14}$ \checkmark \\
C & 孪生增长率 vs Benettin $\lambda$$_1$ & 孪生轨迹的有限时间增长率必须与 Benettin 谱给出的 $\lambda$$_1$ 同量级(\allowbreak{}两者是同一个量的两种测法：一个直接量两条轨道之差，一个走变分方程 + QR) & 0.0236775 \checkmark \\
C & 饱和距离(尾段平均分离) & 孪生轨迹的有限时间增长率必须与 Benettin 谱给出的 $\lambda$$_1$ 同量级(\allowbreak{}两者是同一个量的两种测法：一个直接量两条轨道之差，一个走变分方程 + QR) & 13.8653 \checkmark \\
C & $\lambda$(4) 与 ln2 之差 vs 分块标准误 & 有限时间 $\lambda$ 估计的误差应当落在分块标准误的量级内：块标准误度量的是同一件事(\allowbreak{}有限时间涨落)，如果误差远大于它，说明估计有系统偏差 & 0.219304 \checkmark \\
C & Kaplan--Yorke 维数(Lorenz，由实测谱推出) & 有限时间 $\lambda$ 估计的误差应当落在分块标准误的量级内：块标准误度量的是同一件事(\allowbreak{}有限时间涨落)，如果误差远大于它，说明估计有系统偏差 & 2.06249 \checkmark \\
A & A=QR 的重建误差 & 自写的修正 Gram--Schmidt 必须满足 A=\allowbreak{}QR 与 Q\textasciicircum{}T Q=\allowbreak{}I；det 的绝对值等于 $\Pi$|R\_ii| & $1.110\times 10^{-16}$ \checkmark \\
A & |det A| vs $\Pi$|R\_\allowbreak{}ii| & 自写的修正 Gram--Schmidt 必须满足 A=\allowbreak{}QR 与 Q\textasciicircum{}T Q=\allowbreak{}I；det 的绝对值等于 $\Pi$|R\_ii| & $8.882\times 10^{-16}$ \checkmark \\
A & 数值散度 vs 闭式 $-$13.6667 & 自写的修正 Gram--Schmidt 必须满足 A=\allowbreak{}QR 与 Q\textasciicircum{}T Q=\allowbreak{}I；det 的绝对值等于 $\Pi$|R\_ii| & $1.466\times 10^{-9}$ \checkmark \\
C & 向量化分岔图 vs 标量迭代的最大偏差 & 分岔图的点云与逐参数标量迭代必须给出同一批点(\allowbreak{}向量化路径不能改变结果) & 0 \checkmark \\
\end{longtable}
